\documentclass[12pt]{article}
\usepackage[paperheight=210mm,paperwidth=148mm,top=2cm,bottom=2cm,left=1cm,right=1cm]{geometry}  % A5 paper
\usepackage{background}
\backgroundsetup{
    scale=1,
    angle=0,
    color=black,
    contents={%
    \includegraphics[width=\paperwidth]{../../assets/background}
    }
}
%%%%%%%%%%%%%%%%%%%%%%%%%%%%%%%%%%%%%%%%
\usepackage[utf8]{inputenc}
\usepackage[brazil]{babel}
\usepackage{url}
\usepackage{indentfirst}
\usepackage{titlesec}
% \usepackage[all]{hypcap}
\usepackage{multirow}
% \usepackage{tikz}
\usepackage{ctable}
%%%%%%%%%%%%%%%%%%%%%%%%%%%%%%%%%%%%%%%%
\renewcommand{\familydefault}{\sfdefault}
% \titleformat{\chapter}[display]
% {\bfseries\LARGE\sffamily}{\chaptertitlename}{10pt}{\LARGE}
% \titlespacing*{\chapter}{0pt}{10pt}{10pt}
% \numberwithin{table}{chapter}


\begin{document}
% Cover
\thispagestyle{empty}
\backgroundsetup{
    opacity=1.0,
contents={\includegraphics{../../assets/capa}}
}
\vspace*{12cm}
\begin{Huge}
  Postmortem
\end{Huge}

\newpage

\backgroundsetup{
    opacity=0.1,
contents={
    \includegraphics[width=\paperwidth]{../../assets/background}
    }
    }

\section*{Projeto}

Entre os dias 16 e 20 de maio de 2016 realizou-se a 4a Conferência SciPy Latin America 2016 na Universidade Federal de Santa Catarina. Essa conferência teve como objetivo o fomento da ciência, a adoção da computação como ferramenta para o estudo científico e o uso de Python como principal ferramenta computacional na ciência. Nessa edição contamos com atividades astronômicas para crianças,
um workshop da Software Carpentry,
% FIXME Verificar número de palestras
12 tutoriais e 34 palestras (sendo 18 palestras relâmpago).
% FIXME Incluir o número de horas da trilha para crianças.
Incluindo os eventos sociais tivemos mais de 60h de atividades durante os cinco dias que compuseram o SciPy Latin America 2016.

% FIXME Corrigir número de participantes para incluir as crianças.
Participaram das atividades mais de 70 pessoas dentre
alunos do ensino fundamental,
alunos de graduação,
alunos de pós-graduação,
professores,
desenvolvedores
e entusiastas.

\newpage

\section*{Equipe}

\begin{itemize}
  \item Subcomitê Executivo

    \begin{itemize}
      \item Ivan Ogasawara
      \item Letícia Portella
      \item Raniere Gaia Costa da Silva
    \end{itemize}

  \item Subcomitê Acadêmico

    \begin{itemize}
      \item Antonio Fernando Härter Fetter Filho
    \end{itemize}

  \item Subcomitê Científico

    \begin{itemize}
      \item Melissa Weber Mendonça
    \end{itemize}

  \item Subcomitê de Treinamento

    \begin{itemize}
      \item Jeudy Blanco Vega
    \end{itemize}

  \item Subcomitê Infantil

    \begin{itemize}
      \item Antonio Kanaan
    \end{itemize}

  \item Subcomitê de Comunicação

    \begin{itemize}
      \item Tomas Aliaga
    \end{itemize}

  \item Subcomitê de Apoio Financeiro

    \begin{itemize}
      \item Matheus Braun Magrin
    \end{itemize}

  \item Subcomitê de Suporte Técnico

    \begin{itemize}
      \item Juan Bautista Cabral
    \end{itemize}

\end{itemize}

\newpage

\section*{Infraestrutura}

\subsection*{Sucessos}

\begin{itemize}
  \item Sala para o workshop,

    a infraestrutura do CIASC é muito boa pois inclui mesas confortáveis, várias
    tomadas, ar-condicionado, máquina para café, sanitários próximos.

  \item Auditório,

    o tamanho do auditório, poltronas e equipamentos eletrônicos era adequados.

  \item Servidores de hospedagem web,

    funcionaram adequadamente.

  \item Localização,

    o acesso à salas e auditórios era fácil.
\end{itemize}

\subsection*{Falhas}

\begin{itemize}
  \item Serviço de submissão de palestras,

    que era confuso.

    \textbf{Ações recomendadas:} utilizar um serviço de terceiro como
    \url{http://speakerfight.com/} ou \url{http://www.easychair.org/}.

  \item Sala para os tutoriais,

    mesas eram inadequadas para o apoio dos notebooks e ausência de tomadas.

    \textbf{Ações recomendadas:} Visitar as salas antes antes de reservá-las.

  \item Ausência de cozinha próxima à salas e auditórios,

    o que dificultou o preparo da água para chá, do café e dos alimentos.

    \textbf{Ações recomendadas:} dar preferência para salas e tutoriais próximos
    de uma cozinha.

  \item Número pequeno de bebedouros,

    o que criou longas filas para participantes pegarem água.

    \textbf{Ações recomendadas:} disponibilizar galões de água na ausência de
    bebedouros.

  \item Equipamentos de som,

    salas e auditórios não possuíam equipamento de som para uso fácil pois
    faltaram adaptadores e equipamentos sem fio entravam em conflito.

    \textbf{Ações recomendadas:} Visitar as salas e auditórios antes de
    reservá-los.

  \item Pontos de acesso à internet,

    a conexão estava instável durante vários momentos.

    \textbf{Ações recomendadas:} Testar os pontos de acessos anteriormente.

\end{itemize}

\newpage

\section*{Equipe}

\subsection*{Sucessos}

\begin{itemize}
  \item Organizadores dedicados.
\end{itemize}

\subsection*{Falhas}

\begin{itemize}
  \item Pequeno número de membros.

    \textbf{Ações recomendadas:} Convidar mais pessoas para participar da
    equipe.

  \item Baixa diversidade,

    o número de mulheres envolvidas na organização foi pequeno.

    \textbf{Ações recomendadas:} Convidar mais pessoas para participar da
    equipe.

  \item Comunicação interna,

    lista de tarefas disponíveis era desconhecida e confusa, meios de
    comunicação utilizados eram ineficientes, horários de reuniões online eram
    ruins.

    \textbf{Ações recomendadas:} Escolher melhor as ferramentas utilizadas,
    talvez utilizar kanban.
    E fazer reuniões periódicas.

  \item Ausência de pessoa para controlar o equipamento de som,

    a universidade não forneceu uma pessoa para operar os equipamentos de som.

    \textbf{Ações recomendadas:} requisitar o profissional para a universidade.

  \item Ausência de equipe para filmagem,

    a universidade não forneceu uma equipe e equipamento para filmagem.

    \textbf{Ações recomendadas:} contratar uma equipe externa para essa função.

  \item Ausência de equipe para fotografia,

    nenhum participante tinha câmera profissional.

    \textbf{Ação recomendada:} oferecer vários apoios financeiros para
    participantes voluntários que irão fotografar o evento com suas câmeras
    (semi-)profissionais.

  \item Ausência de pessoa para recepção dos participantes no evento.

    \textbf{Ação recomendada:} oferecer vários apoios financeiros para
    participantes que serão voluntários na recepção do evento.

\end{itemize}

\newpage

\section*{Logística}

% \subsection*{Sucessos}
%
% \begin{itemize}
% \end{itemize}

\subsection*{Falhas}

\begin{itemize}
  \item Utilização dos lugares vagos por não inscritos.

    \textbf{Ações recomendadas:} mencionar no site e outras mídias que é
    possível fazer a inscrição no lugar.

  \item Controle de presença,

    pois as listas de presenças não foram bem divulgadas.

    \textbf{Ações recomendadas:} imprimir as listas de presença com antecedência
    ou utilizar o aplicativo para dispositivos móveis.

  \item Deslocamento dentro da cidade.

    \textbf{Ações recomendadas:} contratar serviço de transporte para os
    participantes puderem ir em lugares mais distantes.

\end{itemize}

\newpage

\section*{Comunicação}

\subsection*{Sucessos}

\begin{itemize}
  \item Vários tweets antes e durante o evento.

  \item Mapas indicando as salas e auditórios no site.
\end{itemize}

\subsection*{Falhas}

\begin{itemize}
  \item Nome da conferência,

    que não remete a nada amplamente conhecido.

    \textbf{Ações recomendadas:} utilizar a marca PyData ao invés do invés de
    SciPy.

  \item Inscrições,

    inscrições para tutoriais e workshops foi confusa.

    \textbf{Ações recomendadas:} simplificar o processo de inscrição.

  \item Informações no site,

    algumas informações como localização de salas e auditórios estava difícil de
    serem encontradas.

    \textbf{Ações recomendadas:} conseguir ajuda de alguém com experiência em
    ``user experience''.

  \item Serviço de chat,

    o aplicativo do Gitter para dispositivos móveis não funcionou adequadamente,
    os participantes não recebiam notificação de novas mensagens e não tinha
    opção para enviar localização, imagens e copiar texto.

    \textbf{Ações recomendadas:} utilizar um serviço cujo aplicativo para
    dispositivas móveis funcione melhor.

  \item Divulgação,

    a divulgação foi fraca fora dos lugares onde já existe usuários.

    \textbf{Ações recomendadas:} uso de cartazes em lugares de grande circulação
    de pessoas e publicações nas redes sociais.

  \item Publicação e edição de conteúdo no site,

    que precisava de conhecimento técnico.

    \textbf{Ações recomendadas:} utilizar um CMS.

  \item Tradução do site,

    que precisava de conhecimento técnico.

    \textbf{Ações recomendadas:} utilizar um serviço de tradução mais amigável.
\end{itemize}

\newpage

\section*{Financeiro}

\subsection*{Sucessos}

\begin{itemize}
  \item Número de patrocinadores,

    esteve dentro do esperado com base nos anos anteriores.
\end{itemize}

\subsection*{Falhas}

\begin{itemize}
  \item Início do envio das cartas de patrocínio,

    pois várias empresas fecham o orçamento para patrocínio no final do ano
    anterior.

    \textbf{Ações recomendadas:} Enviar as cartas de patrocínio no mês de
    Outubro.

  \item Ausência de patrocínio local,

    a maioria dos patrocinadores eram empresas sediadas nos Estados Unidos da
    America.

    \textbf{Ações recomendadas:} Maior divulgação para as empresas locais.

  \item Demora para localização da entrada de divisas,

    por problemas internos da equipe financeira.

    \textbf{Ações recomendadas:} Utilizar uma equipe financeira mais eficiente e
    com conhecimento de formulários estrangeiros como o W8.

  \item Restrições nos pagamentos,

    as regras da equipe financeira eram muito restritivas.

    \textbf{Ações recomendadas:} Utilizar uma equipe financeira cujas regras
    sejam mais simples e flexíveis como ausência da necessidade de três
    orçamentos e pagamento de diárias em valores não tabelados.

\end{itemize}

\newpage

\section*{Outras considerações para próximos eventos}

\begin{itemize}
  \item Cobrar ingresso dos participantes,

    pois isso irá aumentar a taxa de presentes e facilitará a conquista de
    patrocinadores locais.

  \item Oferecer mais benefícios para os patrocinadores,

    como isenção de inscrições e hospedagem fornecida pela organização.

\end{itemize}

\clearpage
\newpage

% Cover
\thispagestyle{empty}
\backgroundsetup{
    opacity=1.0,
contents={\includegraphics{../../assets/contra-capa}}
}
\phantom{1}  % Necessário para imprimir o fundo.

\end{document}
